% 5-appendix-c.tex starts here
% \raggedbottom for the appendix: the intro paragraph shares a page with no
% float (the ER diagram is [p], on its own page), so \flushbottom from the body
% would stretch that short paragraph's vertical glue into big gaps. Previously
% set by the now-removed LaTeX-tooling and PAN-manifest appendices.
\raggedbottom

\section{The data model}\label{anexa:model-date}

The complete database schema --- the 20 \gls{prisma} models --- is shown in
\figref{fig:er-diagram} as an entity-relationship diagram. The diagram is
generated automatically from \texttt{schema.prisma} (with
\textit{prisma-dbml-generator}) and rendered with \textit{dbdocs.io}, so it
stays in sync with the actual data model. Primary keys are marked with the key
symbol, foreign keys with the link symbol, and the connectors between tables
show the relationships between entities; the junction tables (for example
\texttt{UserDietaryRestriction}, \texttt{RecipeCollectionEntry}) are placed next
to the entities they link.

\begin{figure}[p]
  \centering
  % Image is the dbdocs PDF export cropped to its content with `pdfcrop` ---
  % the raw export carries a full-canvas white border that wastes ~1/3 of the
  % page. Re-crop after any re-export. keepaspectratio + the height cap keep
  % the diagram inside the text block.
  \makebox[\textwidth]{%
    \includegraphics[width=\textwidth,height=0.9\textheight,keepaspectratio]{kookio-er-diagram}}
  \caption[Entity-relationship diagram of the database]{Entity-relationship
    diagram of the data schema: the 20 models and the relationships between
    them, generated from \texttt{schema.prisma} and published on
    \textit{dbdocs.io}. The interactive, browsable version:
    \url{https://dbdocs.io/dinea.eduard/kookio-db} (access password:
    \texttt{FMI\_2026}).}
  \label{fig:er-diagram}
\end{figure}

% 5-appendix-c.tex ends here
